% Cache coherence: a write by one CPU makes the cached copies of the same
% memory location on the other CPUs stale; the hardware must update or
% invalidate them.
% Usage: % Cache coherence: a write by one CPU makes the cached copies of the same
% memory location on the other CPUs stale; the hardware must update or
% invalidate them.
% Usage: % Cache coherence: a write by one CPU makes the cached copies of the same
% memory location on the other CPUs stale; the hardware must update or
% invalidate them.
% Usage: % Cache coherence: a write by one CPU makes the cached copies of the same
% memory location on the other CPUs stale; the hardware must update or
% invalidate them.
% Usage: \input{figures/l10-coherence}   (width ~116 mm)
\begin{tikzpicture}[x=1mm, y=1mm, font=\footnotesize\sffamily,
  cpu/.style={draw, fill=white, minimum width=30mm, minimum height=6mm, inner sep=1pt},
  cache/.style={draw, fill=ln-box, minimum width=30mm, minimum height=6mm,
                inner sep=1pt, font=\scriptsize\sffamily},
  lab/.style={font=\scriptsize\sffamily, align=center},
  >={Stealth[length=1.8mm]}]
  \node[cpu] (c1) at (17,14) {CPU 1: \texttt{test\_and\_set}};
  \node[cpu] (c2) at (58,14) {CPU 2};
  \node[cpu] (c3) at (99,14) {CPU 3};
  \node[cache, anchor=north, fill=white, thick] (k1) at (c1.south)
    {\texttt{lock = BUSY}};
  \node[cache, anchor=north] (k2) at (c2.south)
    {\texttt{lock = FREE}\ \iflangja{（古い）}{(stale)}};
  \node[cache, anchor=north] (k3) at (c3.south)
    {\texttt{lock = FREE}\ \iflangja{（古い）}{(stale)}};
  \draw[line width=1.6pt] (0,-6) -- (116,-6);
  \foreach \k in {1,2,3} { \draw (k\k.south) -- (k\k.south |- 0,-6); }
  \node[draw, fill=ln-box, minimum width=40mm, minimum height=6mm] (m) at (58,-15)
    {\iflangja{メモリ}{memory}: \texttt{lock = BUSY}};
  \draw (m.north) -- (58,-6);
  % write path
  \draw[->, thick, ln-accent] ([xshift=-6mm]k1.south) -- ++(0,-12) |- (m.west)
    node[lab, pos=0.25, left, text=ln-accent] {\iflangja{(1) 書き込み}{(1) write}};
  % coherence
  \draw[->, thick, ln-caution, densely dashed] (m.east) -| ([xshift=6mm]k3.south);
  \node[lab, text=ln-caution] at (89,-9.5)
    {\iflangja{(2) 更新または無効化}{(2) update or invalidate}};
  \draw[->, thick, ln-caution, densely dashed] ([xshift=8mm]m.north) -- ([xshift=8mm]k2.south);
\end{tikzpicture}
   (width ~116 mm)
\begin{tikzpicture}[x=1mm, y=1mm, font=\footnotesize\sffamily,
  cpu/.style={draw, fill=white, minimum width=30mm, minimum height=6mm, inner sep=1pt},
  cache/.style={draw, fill=ln-box, minimum width=30mm, minimum height=6mm,
                inner sep=1pt, font=\scriptsize\sffamily},
  lab/.style={font=\scriptsize\sffamily, align=center},
  >={Stealth[length=1.8mm]}]
  \node[cpu] (c1) at (17,14) {CPU 1: \texttt{test\_and\_set}};
  \node[cpu] (c2) at (58,14) {CPU 2};
  \node[cpu] (c3) at (99,14) {CPU 3};
  \node[cache, anchor=north, fill=white, thick] (k1) at (c1.south)
    {\texttt{lock = BUSY}};
  \node[cache, anchor=north] (k2) at (c2.south)
    {\texttt{lock = FREE}\ \iflangja{（古い）}{(stale)}};
  \node[cache, anchor=north] (k3) at (c3.south)
    {\texttt{lock = FREE}\ \iflangja{（古い）}{(stale)}};
  \draw[line width=1.6pt] (0,-6) -- (116,-6);
  \foreach \k in {1,2,3} { \draw (k\k.south) -- (k\k.south |- 0,-6); }
  \node[draw, fill=ln-box, minimum width=40mm, minimum height=6mm] (m) at (58,-15)
    {\iflangja{メモリ}{memory}: \texttt{lock = BUSY}};
  \draw (m.north) -- (58,-6);
  % write path
  \draw[->, thick, ln-accent] ([xshift=-6mm]k1.south) -- ++(0,-12) |- (m.west)
    node[lab, pos=0.25, left, text=ln-accent] {\iflangja{(1) 書き込み}{(1) write}};
  % coherence
  \draw[->, thick, ln-caution, densely dashed] (m.east) -| ([xshift=6mm]k3.south);
  \node[lab, text=ln-caution] at (89,-9.5)
    {\iflangja{(2) 更新または無効化}{(2) update or invalidate}};
  \draw[->, thick, ln-caution, densely dashed] ([xshift=8mm]m.north) -- ([xshift=8mm]k2.south);
\end{tikzpicture}
   (width ~116 mm)
\begin{tikzpicture}[x=1mm, y=1mm, font=\footnotesize\sffamily,
  cpu/.style={draw, fill=white, minimum width=30mm, minimum height=6mm, inner sep=1pt},
  cache/.style={draw, fill=ln-box, minimum width=30mm, minimum height=6mm,
                inner sep=1pt, font=\scriptsize\sffamily},
  lab/.style={font=\scriptsize\sffamily, align=center},
  >={Stealth[length=1.8mm]}]
  \node[cpu] (c1) at (17,14) {CPU 1: \texttt{test\_and\_set}};
  \node[cpu] (c2) at (58,14) {CPU 2};
  \node[cpu] (c3) at (99,14) {CPU 3};
  \node[cache, anchor=north, fill=white, thick] (k1) at (c1.south)
    {\texttt{lock = BUSY}};
  \node[cache, anchor=north] (k2) at (c2.south)
    {\texttt{lock = FREE}\ \iflangja{（古い）}{(stale)}};
  \node[cache, anchor=north] (k3) at (c3.south)
    {\texttt{lock = FREE}\ \iflangja{（古い）}{(stale)}};
  \draw[line width=1.6pt] (0,-6) -- (116,-6);
  \foreach \k in {1,2,3} { \draw (k\k.south) -- (k\k.south |- 0,-6); }
  \node[draw, fill=ln-box, minimum width=40mm, minimum height=6mm] (m) at (58,-15)
    {\iflangja{メモリ}{memory}: \texttt{lock = BUSY}};
  \draw (m.north) -- (58,-6);
  % write path
  \draw[->, thick, ln-accent] ([xshift=-6mm]k1.south) -- ++(0,-12) |- (m.west)
    node[lab, pos=0.25, left, text=ln-accent] {\iflangja{(1) 書き込み}{(1) write}};
  % coherence
  \draw[->, thick, ln-caution, densely dashed] (m.east) -| ([xshift=6mm]k3.south);
  \node[lab, text=ln-caution] at (89,-9.5)
    {\iflangja{(2) 更新または無効化}{(2) update or invalidate}};
  \draw[->, thick, ln-caution, densely dashed] ([xshift=8mm]m.north) -- ([xshift=8mm]k2.south);
\end{tikzpicture}
   (width ~116 mm)
\begin{tikzpicture}[x=1mm, y=1mm, font=\footnotesize\sffamily,
  cpu/.style={draw, fill=white, minimum width=30mm, minimum height=6mm, inner sep=1pt},
  cache/.style={draw, fill=ln-box, minimum width=30mm, minimum height=6mm,
                inner sep=1pt, font=\scriptsize\sffamily},
  lab/.style={font=\scriptsize\sffamily, align=center},
  >={Stealth[length=1.8mm]}]
  \node[cpu] (c1) at (17,14) {CPU 1: \texttt{test\_and\_set}};
  \node[cpu] (c2) at (58,14) {CPU 2};
  \node[cpu] (c3) at (99,14) {CPU 3};
  \node[cache, anchor=north, fill=white, thick] (k1) at (c1.south)
    {\texttt{lock = BUSY}};
  \node[cache, anchor=north] (k2) at (c2.south)
    {\texttt{lock = FREE}\ \iflangja{（古い）}{(stale)}};
  \node[cache, anchor=north] (k3) at (c3.south)
    {\texttt{lock = FREE}\ \iflangja{（古い）}{(stale)}};
  \draw[line width=1.6pt] (0,-6) -- (116,-6);
  \foreach \k in {1,2,3} { \draw (k\k.south) -- (k\k.south |- 0,-6); }
  \node[draw, fill=ln-box, minimum width=40mm, minimum height=6mm] (m) at (58,-15)
    {\iflangja{メモリ}{memory}: \texttt{lock = BUSY}};
  \draw (m.north) -- (58,-6);
  % write path
  \draw[->, thick, ln-accent] ([xshift=-6mm]k1.south) -- ++(0,-12) |- (m.west)
    node[lab, pos=0.25, left, text=ln-accent] {\iflangja{(1) 書き込み}{(1) write}};
  % coherence
  \draw[->, thick, ln-caution, densely dashed] (m.east) -| ([xshift=6mm]k3.south);
  \node[lab, text=ln-caution] at (89,-9.5)
    {\iflangja{(2) 更新または無効化}{(2) update or invalidate}};
  \draw[->, thick, ln-caution, densely dashed] ([xshift=8mm]m.north) -- ([xshift=8mm]k2.south);
\end{tikzpicture}
